% ======================================================================
% 00-abstract.tex
% ======================================================================
\begin{abstract}
\noindent
Research code usually becomes reusable by accident, if at all. We report a
different path: a groundwater-remediation project in which the code was made
reusable \emph{by construction}, one failure at a time, with an AI coding agent
(the Pi harness running DeepSeek models) doing most of the implementation,
testing and record-keeping. The scientific stack is a verified
advection--dispersion sandbox, a batched finite-volume Darcy/transport solver, a
chance-constrained pump-and-treat optimizer, and 1D/2D/3D experiments.
The reusable artifacts are less glamorous: one NumPy reference implementation
with a parity-tested torch backend, configuration manifests that record seed,
device, dtype and git revision, gate suites that turn physical assumptions into
go/no-go conditions, a seed-sweep harness that reports medians with their spread,
and an artifact-first remote runner that keeps the results of a run that dies.
We show that these mechanisms changed conclusions, not just style: a
single-seed ``32\% improvement'' became a negative result, an FNO that scored
worse than predicting zeros became a 1.5\% operator on one seed (3\% median), a policy that passed its
own reliability gate failed on the independent NumPy reference, and a 3D paper
preset reproduced on a third machine with a different BLAS. We also report where
the agent was confidently wrong, and argue that the useful division of labour is
not ``human writes, AI types'' but \emph{human owns the gates, agent owns the
iteration count, and the repository owns the memory}.
\end{abstract}

\vspace{0.5em}
\noindent\textbf{Keywords:} research software; reproducibility; physics-informed
neural networks; Fourier neural operators; stochastic optimization;
AI-assisted development.
